\setcounter{section}{67}
\section{Equipos personales}

Nombre completo: Entorno de desarrollo PHP

\subsection{Introducción}
\texttt{PHP Hypertext Preprocessor} es un lenguaje de código abierto muy adaptado al desarrollo web. Es un lenguaje interpretado con una sintáxis parecida a C, Java y Pearl y usa tipado dinámico. Es compatible con multitud de bases de datos y multiplataforma. 

\subsection{Lenguaje PHP}
    Es un lenguaje de scripting desde el lado del servidor. Este procesa peticiones, genera contenido dinámico y envia la salida al cliente en forma de html. De esta forma el código fuente queda oculto al cliente. PHP también puede crear imágenes, PDFs, XML o JSON.

    PHP se puede escribir dentro de html con la etiqueta \texttt{<?PHP} y \texttt{?>}. Esto será interpretado y ejecutado por el servidor antes de ser enviado.

\subsection{Arquitectura de aplicaciones PHP}
    Existen varios modelos y arquitecturas que integran PHP. Los mas importantes son:
    \begin{description}
        \item[Cliente SErvidor] El cliente envía solicitudes y php las procesa y responde
        \item[LAMP] Linx Apache MySQL PHP
        \item[Arquitectura] MVC, Componentes o Microservicios
        \item[Sesiones y estado] php es capaz de manejar esto por encima de html.
        \item[Despliegues y escalabilidad] Mediante servidores web, Cachés y Balanceo de carga
    \end{description}

\subsection{Integración con bases de datos}
    PHP permite integrar varias bases de datos:
    \begin{description}
        \item[mysqli] MySQL Improved. Interfaz para el manejo de bases de datos
        \item[PDO] PHP data objects, permite mapear bases de datos con objetos.
        \item[Consultas preparadas]
    \end{description}

\subsection{Evolución y mejoras en PHP 7 y 8}
    \subsubsection{PHP 7}
        Rendimiento mejorado, con un motor de ejecución nuevo. También incluye el spaceship operator <=> (sobre el cual tengo opiniones). Tipos escalares (oh la lá), manejo de errores, operador  nulo y gestión de memoria.
    \subsubsection{PHP 8}
        JIT compilation, atributos, unón de tipos, operador nulo seguro, argumentos nombrados, match (switch con cosas) y mejora de expresiones regulares.

\subsection{Frameworks PHP}
    \subsubsection{Laravel}
        De los mas populares y utilizados, es código abierto. Facilita la integración con bases de datos mapeando objetos con relaciones.

        Maneja rutas de manera sencilla y flexible. Tiene sistemas de autenticación y autorización. También puede manejar Queues y Jobs para trabajos en segundo plano.


    \subsubsection{Symfony}
        Es conocido por su estabilidad y ser altamente configurable. Está basado en MVC, separando lógica de negocio, de servidor, y de la aplicación web.
    \subsubsection{Codeigniter}
        Alto rendimineto y fácil de usar. Es simple y rápido.

\subsection{CMS en PHP}
    Hay multitud de sistemas de gestión de contenido basados en php, listas para gestionar páginas web sin necesidad de conocimientos avanzados de programación.

    \begin{description}
        \item[Wordpress] Es el mas popular del mundo, originalmente para blogs, con gran cantidad de plugins.
        \item[Jooma] Ampliamente utilizado, es algo mas complejo que wordpress, con características para integrar muiltilenguajes, plantillas, extensiones\dots
        \item[Drupal] Ideal para tener un amplio grado de personalización, gran cantidad de módulos
        \item[Magento] Plataforma de comercio electrónico con un alto grado de personalización.
        \item[PrestaShop] Otra plataforma de comercio electrónico.
    \end{description}

\subsection{PHP en aplicaciones escalables}
    Permite crear microservicios independientes con todos los protocolos de despliegue que ello implica. 

    También existen APIs REST con frameworks como LAravel o Symfony. Con múltiples estrategias de autenticación. 

\subsection{Buenas prácticas}
    Estndares PSR son recomendaciones adaptadas por la comunidad para garantizar consistencia y calidad de cdigo. Ver \href{https://www.php-fig.org/psr/}{página de estándares}.

    La gestión de dependencias se hace con Composer, similar a maven.

    EL testing se realiza con PHPUnit.